\chapter{The Multi-Cycle Datapath}

The single-cycle processor wastes time because every instruction takes as
long as the slowest. The multi-cycle approach breaks each instruction into
steps, each taking one \emph{short} clock cycle; different instructions use
different numbers of cycles.

\section{Key Ideas}

\begin{itemize}
  \item Choose the clock period to fit the slowest \emph{step} (roughly one
    major unit: a memory access, a register file access, or an ALU
    operation), not the slowest instruction.
  \item A functional unit can be reused in different cycles of the same
    instruction. So we need only \textbf{one memory} (for both instructions
    and data) and \textbf{one ALU} (it also computes $PC+4$ and the branch
    target).
  \item Values produced in one cycle and needed in a later one must be
    saved in new \textbf{internal registers} at the end of the cycle:
    \begin{itemize}
      \item \textbf{IR} (instruction register) and \textbf{MDR} (memory data register),
      \item \textbf{A} and \textbf{B} (register file outputs),
      \item \textbf{ALUOut} (ALU result).
    \end{itemize}
  \item Extra muxes: the memory address comes from PC or ALUOut
    (\code{IorD}); ALU input A is PC or register A (\code{ALUSrcA}); ALU
    input B is register B, 4, the sign-extended immediate, or the shifted
    immediate (\code{ALUSrcB}, 2 bits); the new PC comes from ALU, ALUOut or
    the jump address (\code{PCSource}).
  \item Control becomes a \textbf{finite state machine}.
\end{itemize}

\begin{figure}[H]
\centering
\begin{tikzpicture}[thick, >=Stealth, font=\small,
  blk/.style={draw, minimum height=1.4cm, minimum width=1.2cm, align=center},
  reg/.style={draw, fill=gray!20, minimum height=1cm, minimum width=0.45cm, font=\scriptsize}]
  \node[blk, minimum width=0.6cm] (pc) at (0,0) {PC};
  \node[blk, minimum height=2cm] (mem) at (1.8,0) {Memory};
  \node[reg] (ir) at (3.4,0.5) {IR};
  \node[reg] (mdr) at (3.4,-0.8) {MDR};
  \node[blk, minimum height=2cm] (rf) at (5.2,0) {Regs};
  \node[reg] (a) at (6.6,0.5) {A};
  \node[reg] (b) at (6.6,-0.5) {B};
  \node[draw, minimum height=1.6cm, minimum width=0.9cm] (alu) at (8.2,0) {ALU};
  \node[reg, minimum width=0.9cm] (ao) at (9.8,0) {ALUOut};
  \draw[->] (pc) -- (mem);
  \draw[->] ($(mem.east)+(0,0.5)$) -- (ir);
  \draw[->] ($(mem.east)+(0,-0.8)$) -- (mdr);
  \draw[->] (ir) -- ($(rf.west)+(0,0.5)$);
  \draw[->] (mdr) -- ($(rf.west)+(0,-0.8)$);
  \draw[->] ($(rf.east)+(0,0.5)$) -- (a);
  \draw[->] ($(rf.east)+(0,-0.5)$) -- (b);
  \draw[->] (a) -- ($(alu.west)+(0,0.5)$);
  \draw[->] (b) -- ($(alu.west)+(0,-0.5)$);
  \draw[->] (alu) -- (ao);
  \draw[->] (ao.east) -- ++(0.4,0) |- ++(0,-1.8) -| (mem.south);
\end{tikzpicture}
\caption{Sketch of the multi-cycle datapath: one memory, one ALU, and
internal registers (shaded) between the stages. Muxes are omitted.}
\end{figure}

\section{The Five Steps}

\begin{center}
\begin{tabular}{p{2.2cm}p{3.1cm}p{3.1cm}p{2.6cm}p{2.2cm}}
\toprule
Step & R-type & \code{lw}/\code{sw} & \code{beq} & \code{j} \\
\midrule
1. Fetch & \multicolumn{4}{p{11cm}}{IR $\gets$ Mem[PC]; PC $\gets$ PC + 4} \\
2. Decode & \multicolumn{4}{p{11cm}}{A $\gets$ Reg[rs]; B $\gets$ Reg[rt];
  ALUOut $\gets$ PC + (signext(imm) $\ll$ 2)} \\
3. Execute & ALUOut $\gets$ A op B & ALUOut $\gets$ A + signext(imm) &
  if (A == B) PC $\gets$ ALUOut & PC $\gets$ \{PC[31:28], addr, 00\} \\
4. Memory / R-completion & Reg[rd] $\gets$ ALUOut &
  \code{lw}: MDR $\gets$ Mem[ALUOut]; \code{sw}: Mem[ALUOut] $\gets$ B & & \\
5. Write-back & & \code{lw}: Reg[rt] $\gets$ MDR & & \\
\bottomrule
\end{tabular}
\end{center}

Cycle counts: \code{lw} 5, \code{sw} 4, R-type 4, \code{beq} 3, \code{j} 3.

\begin{remark}
In the decode step we compute the branch target \emph{speculatively}: the
ALU is idle anyway, and we do not yet know whether the instruction is a
branch. If it is not, the result is simply ignored. This is an early
example of doing work optimistically.
\end{remark}

\section{Control as an FSM}

The control unit is a Moore FSM whose states correspond to the steps above.
States 0 (fetch) and 1 (decode) are shared by all instructions; from state 1
the next state depends on the opcode:

\begin{figure}[H]
\centering
\begin{tikzpicture}[->, >=Stealth, thick, node distance=1.6cm and 2.2cm,
  every state/.style={minimum size=1cm, font=\scriptsize, align=center}]
  \node[state, initial] (s0) {0\\Fetch};
  \node[state, right=of s0] (s1) {1\\Decode};
  \node[state, below left=of s1] (s2) {2\\MemAddr};
  \node[state, below=of s1] (s6) {6\\Exec};
  \node[state, below right=of s1] (s8) {8\\Branch};
  \node[state, right=of s8] (s9) {9\\Jump};
  \node[state, below left=of s2] (s3) {3\\MemRd};
  \node[state, below=of s2] (s5) {5\\MemWr};
  \node[state, below=of s3] (s4) {4\\WB};
  \node[state, below=of s6] (s7) {7\\R-done};
  \path (s0) edge (s1)
        (s1) edge node[left, font=\scriptsize]{lw/sw} (s2)
        (s1) edge node[right, font=\scriptsize]{R} (s6)
        (s1) edge node[right, font=\scriptsize]{beq} (s8)
        (s1) edge[bend left] node[above, font=\scriptsize]{j} (s9)
        (s2) edge node[left, font=\scriptsize]{lw} (s3)
        (s2) edge node[right, font=\scriptsize]{sw} (s5)
        (s3) edge (s4)
        (s6) edge (s7);
\end{tikzpicture}
\caption{Multi-cycle control FSM. States 4, 5, 7, 8 and 9 all return to
state 0 (arrows omitted).}
\end{figure}

\begin{center}
\begin{tabular}{cp{13cm}}
\toprule
State & Asserted signals (main ones) \\
\midrule
0 & MemRead, IRWrite, IorD=0, ALUSrcA=0, ALUSrcB=01, ALUOp=00, PCWrite, PCSource=00 \\
1 & ALUSrcA=0, ALUSrcB=11, ALUOp=00 \\
2 & ALUSrcA=1, ALUSrcB=10, ALUOp=00 \\
3 & MemRead, IorD=1 \\
4 & RegWrite, MemtoReg=1, RegDst=0 \\
5 & MemWrite, IorD=1 \\
6 & ALUSrcA=1, ALUSrcB=00, ALUOp=10 \\
7 & RegWrite, MemtoReg=0, RegDst=1 \\
8 & ALUSrcA=1, ALUSrcB=00, ALUOp=01, PCWriteCond, PCSource=01 \\
9 & PCWrite, PCSource=10 \\
\bottomrule
\end{tabular}
\end{center}

The PC is written when $\code{PCWrite} + (\code{PCWriteCond} \cdot
\code{Zero})$.

\paragraph{Implementing the FSM.} Either as hardwired logic (a state
register plus combinational next-state/output logic, e.g.\ a PLA or ROM),
or by \textbf{microprogramming}: each state is a microinstruction stored in
a control ROM, and a micro-PC sequences through them with dispatch tables
for opcode-based branching. Microprogramming was popular for complex CISC
instruction sets.

\section{Performance Comparison}

\begin{example}
Instruction mix: 25\% \code{lw}, 10\% \code{sw}, 45\% R-type, 15\% \code{beq},
5\% \code{j}. CPI $= 0.25(5) + 0.10(4) + 0.45(4) + 0.15(3) + 0.05(3)
= 1.25 + 0.4 + 1.8 + 0.45 + 0.15 = 4.05$.

If the multi-cycle clock is 200\,ps (one memory access) and the
single-cycle clock is 800\,ps:
\begin{itemize}
  \item single-cycle: $1 \times 800 = 800$\,ps per instruction;
  \item multi-cycle: $4.05 \times 200 = 810$\,ps per instruction.
\end{itemize}
Not a win here! The multi-cycle design pays overhead (register setup,
uneven step lengths). Its real advantages are the hardware savings and the
ability to have instructions of very different lengths (e.g.\ floating
point divide) without slowing everything else.
\end{example}

\begin{important}[Exam tip]
To compare designs, always compute time per instruction (or total time) =
CPI $\times$ cycle time, never CPI alone or clock rate alone.
\end{important}

\section{Exceptions (brief)}

An \textbf{exception} (internal, e.g.\ overflow or undefined opcode) or
\textbf{interrupt} (external, e.g.\ I/O) transfers control to the OS. MIPS
saves the faulting PC in the \textbf{EPC} register, records the cause in
the \textbf{Cause} register and jumps to a fixed handler address
(0x80000180). In the multi-cycle FSM this adds states that set EPC, Cause
and the PC.

\section*{Summary}
\begin{itemize}
  \item Multi-cycle: short clock (one major unit per cycle), one memory, one
    ALU, internal registers IR, MDR, A, B, ALUOut.
  \item Steps: fetch, decode/register fetch, execute, memory, write-back;
    \code{lw} 5 cycles, R-type and \code{sw} 4, \code{beq} and \code{j} 3.
  \item Control is an FSM (hardwired or microprogrammed).
  \item Average CPI $> 1$; whether it beats single-cycle depends on the mix
    and the overhead. Pipelining is the real improvement.
\end{itemize}
